\documentclass[11pt]{scrartcl}
\usepackage{napkin-style}
\usepackage{napkin-macros}

\title{Napkin-Style \LaTeX{} Template: A Guide}
\author{}
\date{\shortdate}

\begin{document}
\maketitle
\tableofcontents
\newpage

% ============================================================
\section{Overview}
% ============================================================

This template is adapted from Evan Chen's \emph{An Infinitely Large Napkin}
(\url{https://web.evanchen.cc/napkin.html}), a beautifully typeset open-access
mathematics textbook. The template provides:

\begin{itemize}
    \item Colored theorem, lemma, and example boxes using \texttt{tcolorbox}.
    \item Purple \S~section numbering in the KOMA-Script style.
    \item A rich collection of math macros for algebra, topology, and more.
    \item Sensible hyperlink colors and cross-referencing via \texttt{cleveref}.
    \item Support for commutative diagrams via \texttt{tikz-cd}.
\end{itemize}

% ============================================================
\section{File Structure}
% ============================================================

The template consists of four files:

\begin{center}
\begin{tabular}{ll}
    \textbf{File} & \textbf{Purpose} \\
    \hline
    \texttt{napkin-style.sty} & Visual style: layout, colors, theorem boxes \\
    \texttt{napkin-macros.sty} & Math macros and operator definitions \\
    \texttt{template.tex} & Minimal working document to start from \\
    \texttt{guide.tex} & This documentation file \\
\end{tabular}
\end{center}

To start a new document, copy \texttt{template.tex} to your project folder.
Both \texttt{.sty} files must be in the same folder as your \texttt{.tex} file,
or installed to your local \texttt{texmf} tree (see \Cref{sec:install}).

% ============================================================
\section{How to Compile}
\label{sec:compile}
% ============================================================

The recommended way to compile any document using this template is with
\texttt{latexmk}. Create a \texttt{.latexmkrc} file in your project directory:

\begin{verbatim}
push @extra_pdflatex_options, '-synctex=1';
$pdf_mode = 1;
$max_repeat = 5;
\end{verbatim}

Then run:
\begin{verbatim}
latexmk -f yourfile.tex
\end{verbatim}

\texttt{latexmk} automatically reruns \texttt{pdflatex} as many times as
needed to resolve all cross-references and the table of contents.

\subsection{Manual compilation}

If you prefer to call \texttt{pdflatex} directly, run it \emph{twice}:
\begin{verbatim}
pdflatex yourfile.tex
pdflatex yourfile.tex
\end{verbatim}
The first pass generates auxiliary files; the second pass reads them back
to produce correct page numbers and hyperlinks.

\subsection{Installing the style files globally}
\label{sec:install}

To use the style files from any directory without copying them each time,
place them in your local \texttt{texmf} tree:
\begin{verbatim}
mkdir -p ~/texmf/tex/latex/napkin-style
cp napkin-style.sty napkin-macros.sty ~/texmf/tex/latex/napkin-style/
texhash ~/texmf
\end{verbatim}
After this, you can \verb|\usepackage{napkin-style}| from any project.

% ============================================================
\section{Document Classes}
% ============================================================

The style works with both KOMA-Script classes.

\begin{center}
\begin{tabular}{lll}
    \textbf{Class} & \textbf{Best for} & \textbf{Header behavior} \\
    \hline
    \texttt{scrartcl} & Problem sets, short notes, single-topic & Section title \\
    \texttt{scrbook}  & Multi-chapter notes, course notes & Chapter + section \\
\end{tabular}
\end{center}

The style file detects the class automatically. For a book-style document:
\begin{verbatim}
\documentclass[11pt,twoside=semi,openright]{scrbook}
\usepackage{napkin-style}
\usepackage{napkin-macros}
\end{verbatim}

% ============================================================
\section{Theorem Environments}
% ============================================================

All environments are numbered together within each section, so the counter
resets at each new section. The available environments are listed below with
live examples.

\subsection{Blue theorem box (theorem, lemma, proposition, corollary)}

\begin{theorem}[Lagrange]
    If $G$ is a finite group and $H \le G$, then $\abs{H}$ divides $\abs{G}$.
\end{theorem}

\begin{lemma}
    Every group of prime order is cyclic.
\end{lemma}

\begin{proposition}
    The center $Z(G)$ is always a normal subgroup of $G$.
\end{proposition}

\begin{corollary}
    If $\abs{G} = p$ is prime, then $G \cong \Zc{p}$.
\end{corollary}

\subsection{Orange example box}

\begin{example}
    The symmetric group $S_3$ has order $6 = 3!$. Its center is trivial:
    $Z(S_3) = \{e\}$.
\end{example}

\subsection{Green remark box}

\begin{remark}
    Lagrange's theorem has no converse in general: a group of order $6$
    need not have a subgroup of every order dividing $6$.
\end{remark}

\subsection{Question and exercise box}

\begin{ques}
    Does every group of order $4$ have a subgroup of order $2$?
\end{ques}

\begin{exercise}
    Prove that any group of order $p^2$ (for prime $p$) is abelian.
\end{exercise}

\subsection{Definition and claim (plain style)}

\begin{definition}
    A \vocab{normal subgroup} $N \normalin G$ satisfies $gNg\inv = N$
    for all $g \in G$.
\end{definition}

\begin{claim}
    The kernel of any homomorphism is normal.
\end{claim}

\subsection{Subproof}

Inside a proof, you can use \texttt{subproof} to delimit a sub-argument
with a filled square $\blacksquare$ as the QED symbol:

\begin{proof}
    We show two things.

    \begin{subproof}[Step 1]
        The kernel is a subgroup (closure under product and inverse).
    \end{subproof}

    \begin{subproof}[Step 2]
        It is closed under conjugation because homomorphisms respect products.
    \end{subproof}

    Both steps together complete the proof.
\end{proof}

% ============================================================
\section{Special Commands}
% ============================================================

\subsection{\texttt{\textbackslash prototype}}

Use \verb|\prototype{...}| at the start of a section to highlight the
canonical motivating example. It renders in red italic.

\prototype{The integers $\ZZ$ under addition.}

The text of the section then continues normally after the prototype line.

\subsection{\texttt{moral} environment}

The \texttt{moral} environment draws a green bordered box for the key
conceptual takeaway of a section.

\begin{moral}
    Normal subgroups are exactly the subgroups that arise as kernels
    of homomorphisms.
\end{moral}

\subsection{\texttt{\textbackslash vocab}}

\verb|\vocab{term}| typesets a vocabulary term in \vocab{blue bold},
useful when introducing a new definition inline.

\subsection{\texttt{\textbackslash hrulebar}}

\verb|\hrulebar| inserts a thin horizontal rule, useful as a visual separator.

\hrulebar

Text continues here after the rule.

% ============================================================
\section{Commutative Diagrams}
% ============================================================

The \texttt{tikz-cd} package is loaded and configured with hook arrows
matching Napkin's style.

\begin{example}
    The first isomorphism theorem gives the commutative diagram
    \[
    \begin{tikzcd}
        G \arrow[r, "\phi"] \arrow[d, "\pi"', two heads]
            & H \\
        G/\ker\phi \arrow[ru, "\bar\phi"', hook]
    \end{tikzcd}
    \]
    where $\pi$ is the quotient map and $\bar\phi$ is the induced injection.
\end{example}

Useful \texttt{tikz-cd} arrow options:

\begin{center}
\begin{tabular}{ll}
    \verb|two heads|    & surjective $\twoheadrightarrow$ \\
    \verb|hook|         & injective $\hookrightarrow$ \\
    \verb|dashed|       & dashed arrow \\
    \verb|squiggly|     & squiggly arrow \\
    \verb|Rightarrow|   & double arrow $\Rightarrow$ \\
\end{tabular}
\end{center}

% ============================================================
\section{Math Macros Reference}
% ============================================================

\subsection{Number fields}

\begin{center}
\begin{tabular}{lll}
    \textbf{Command} & \textbf{Output} & \textbf{Meaning} \\
    \hline
    \verb|\NN| & $\NN$ & Natural numbers \\
    \verb|\ZZ| & $\ZZ$ & Integers \\
    \verb|\QQ| & $\QQ$ & Rationals \\
    \verb|\RR| & $\RR$ & Reals \\
    \verb|\CC| & $\CC$ & Complex numbers \\
    \verb|\FF| & $\FF$ & Generic field \\
    \verb|\EE| & $\EE$ & Expectation / generic field \\
\end{tabular}
\end{center}

\subsection{Common operators}

\begin{center}
\begin{tabular}{lll}
    \verb|\abs{x}|    & $\abs{x}$    & Absolute value \\
    \verb|\norm{x}|   & $\norm{x}$   & Norm \\
    \verb|\floor{x}|  & $\floor{x}$  & Floor \\
    \verb|\ceiling{x}|& $\ceiling{x}$& Ceiling \\
    \verb|\cbrt{x}|   & $\cbrt{x}$   & Cube root \\
    \verb|\eps|       & $\eps$       & $\varepsilon$ \\
    \verb|\defeq|     & $\defeq$     & Definition equality \\
    \verb|\injto|     & $\injto$     & Injection arrow \\
    \verb|\surjto|    & $\surjto$    & Surjection arrow \\
\end{tabular}
\end{center}

\subsection{Algebra shorthands}

\begin{center}
\begin{tabular}{lll}
    \verb|\Zc{n}|     & $\Zc{n}$     & $\ZZ/n\ZZ$ \\
    \verb|\Zm{n}|     & $\Zm{n}$     & $(\ZZ/n\ZZ)^\times$ \\
    \verb|\normalin|  & $\normalin$  & Normal subgroup \\
    \verb|\GL|, \verb|\SL|, \verb|\Gal|, \verb|\Aut|, \verb|\Hom|, \verb|\End|
                      & \multicolumn{2}{l}{Standard algebra operators} \\
\end{tabular}
\end{center}

\subsection{Category theory}

\begin{center}
\begin{tabular}{lll}
    \verb|\catname{Grp}| & $\catname{Grp}$ & Category name (sans-serif) \\
    \verb|\op|           & $\cdot^{\op}$   & Opposite category \\
    \verb|\obj|          & $\obj$          & Object class \\
    \verb|\Hom|          & $\Hom$          & Hom-set \\
\end{tabular}
\end{center}

\subsection{Script letters}

\verb|\SA|, \verb|\SB|, \verb|\SC|, \verb|\SD|, \verb|\SF|, \verb|\SG|, \verb|\SH|
produce $\SA, \SB, \SC, \SD, \SF, \SG, \SH$ (mathscr).

\verb|\AA|, \verb|\BB|, \verb|\OO|, \verb|\PP|, \verb|\MM|
produce $\AA, \BB, \OO, \PP, \MM$ (mathcal).

\subsection{Cross-referencing with \texttt{cleveref}}

Instead of \verb|\ref|, use \verb|\cref| and \verb|\Cref|:

\begin{center}
\begin{tabular}{ll}
    \verb|\cref{thm:lagrange}|  & ``theorem~3.1'' (lowercase) \\
    \verb|\Cref{thm:lagrange}|  & ``Theorem~3.1'' (capitalized) \\
\end{tabular}
\end{center}

\texttt{cleveref} automatically inserts the correct environment name.
Always label your environments:
\begin{verbatim}
\begin{theorem}
\label{thm:my-theorem}
    ...
\end{theorem}
\end{verbatim}

% ============================================================
\section{Customization}
% ============================================================

\subsection{Changing theorem colors}

The box styles are defined in \texttt{napkin-style.sty} as \texttt{tcbset}
keys. To change the theorem box color, add this after \verb|\usepackage{napkin-style}|:
\begin{verbatim}
\tcbset{
    theorem box/.style={
        enhanced, arc=9pt, colframe=Plum, colback=Plum!5,
        boxrule=1pt, before skip=12pt, after skip=14pt,
        left=10pt, right=10pt
    }
}
\end{verbatim}

\subsection{Changing the \S~symbol}

To remove the section symbol prefix, comment out or change these lines
in \texttt{napkin-style.sty}:
\begin{verbatim}
\renewcommand*{\sectionformat}{\color{purple}\S\thesection\autodot\enskip}
\renewcommand*{\subsectionformat}{\color{purple}\S\thesubsection\autodot\enskip}
\end{verbatim}

\subsection{Adding a bibliography}

To add references, load \texttt{biblatex} after both style packages:
\begin{verbatim}
\usepackage{napkin-style}
\usepackage{napkin-macros}
\usepackage[backend=biber, style=alphabetic]{biblatex}
\addbibresource{references.bib}
\end{verbatim}
Then place \verb|\printbibliography| at the end of the document, and
compile with:
\begin{verbatim}
latexmk -f yourfile.tex   % latexmk runs biber automatically
\end{verbatim}

\end{document}
